% =====================================================
\chapter{Literature Survey}
\label{chap:literature_survey}


This chapter builds the theoretical foundation for the work presented in this report.  

% =====================================================
\section{Warehouse Digital Twins}
\label{sec:Warehouse DT}

The first paper reviewed, authored by Ferrari et al. \cite{Ferrari2022}, has focused on the creation 
of a Digital Shadow for an Automated Storage and Retrieval System (AS/RS) in a university
logistics lab. A detailed simulation model was developed using AnyLogic, incorporating 
discrete event and agent-based modeling. Although real-time control was not achieved, the 
model was able to provide insights into system layout and performance. A Warehouse Management 
System (WMS) and a Warehouse Control System (WCS) were used to coordinate operations such as 
order generation and storage location management. The system was designed to simulate the behavior 
of physical operations closely, though without live data feedback into the physical system. 

The approach by Lugaresi et al. \cite{Lugaresi2019}, involved a small-scale manufacturing setup using LEGO components.
A closed-loop production system with four stations (S1, S2, S3, S4) and associated buffers was designed,
controlled by LEGO® MINDSTORMS® EV3® bricks programmed with Python scripts. Two product types (A and B) 
with distinct processing times, modeled using Weibull distributions, were processed, and setup times for 
switching between product types were required by stations S2, S3, and S4, with S2’s setup time increasing 
exponentially after three consecutive type A pallets. A discrete event simulator, developed in ARENA® (Sim.doe), 
was used as the digital model, synchronized with the physical LMS through data exchange files (*.txt) containing 
buffer levels, processing times, and production status. Real-time data from the LMS were collected, the digital 
model was updated, and alternative production policies were simulated to identify the optimal response to a disruption
in the RTS framework.

In the work by Zhang et al, \cite{Zhang2026} the core idea is \emph{opti-state control}. Rather than a single optimization before execution, the system is continuously monitored.
 As the warehouse is deviated from its desired performance, there is adaptive reconfiguration of plans.  For warehouse and intralogistics research, the paper's contribution is to 
 move from digital twins as a monitoring tool, toward a closed-loop control system.
% =====================================================

% =====================================================
\section*{Summary}

Research has steadily advanced toward controllable and interpretable generation.  
Transformers introduced attention mechanisms, ViTs extended them to vision, and GANs and VAEs pioneered generative modeling.  
Diffusion unified these ideas into a stable probabilistic framework, enhanced by latent conditioning, cross-attention, and adapters like ControlNet and IP-Adapter.  
Contrastive systems such as CLIP and CSD added semantic and stylistic grounding, forming the basis for the adaptive, CSD-guided diffusion explored here.
